\section{The Lean Development}
\label{app:lean}

The development is a Lean~4 library over Mathlib, with the root namespace
\lean{DeciSigma1}; all names below are relative to it. \Cref{tab:lean-layout}
lists its parts, \cref{tab:lean-defs,tab:lean-thms} relate the definitions
and results of the paper to declarations, and \cref{app:lean-differences}
lists the differences between the two presentations.

{\small
\begin{longtable}{@{}>{\raggedright\arraybackslash}p{0.22\linewidth}>{\raggedright\arraybackslash}p{0.74\linewidth}@{}}
\caption{The parts of the development.}\label{tab:lean-layout}\\
\toprule
Part & Contents \\
\midrule
\endfirsthead
\toprule
Part & Contents \\
\midrule
\endhead
\bottomrule
\endfoot
\lean{Main.lean} & the public statements, with a map of the proof \\*
\lean{Encodings.lean} & satisfiability and entailment as sentences \\*
\lean{Semantics} & trees over both signatures, the order with variance, finite and regular trees \\*
\lean{Constraints} & flat literals, closure, the signed translation, compiled bounds by finite trees \\*
\lean{Satisfiability} & satisfiability and finite satisfiability \\*
\lean{Automata} & runs of transition functions, the run automata of constraints \\*
\lean{Words} & periods, primitive words, rays, deletion of positions, suffix chains \\*
\lean{Arithmetic} & difference constraints, simple paths, maximal repetitions, iterated exponentials \\*
\lean{Readers} & tables, structural readers, holes, domination \\*
\lean{Holes} & the small-hole theorem: \lean{Fresh}, \lean{Periodic}, \lean{Tail}, and the bounds \\*
\lean{Full} & the reduction to holes over $\BFT$, regular and finite countermodels, deciders \\*
\lean{Entailment} & the signature $\BF$ and its reduction by a guard \\*
\lean{Logic} & sentences, normal forms, the decision procedure for sentences \\*
\lean{Examples}, \lean{Audit.lean} & examples; the check of axioms and of computability \\
\end{longtable}}

{\small
\begin{longtable}{@{}>{\raggedright\arraybackslash}p{0.30\linewidth}>{\raggedright\arraybackslash\footnotesize}p{0.66\linewidth}@{}}
\caption{Definitions.}\label{tab:lean-defs}\\
\toprule
Paper & \small Lean \\
\midrule
\endfirsthead
\toprule
Paper & \small Lean \\
\midrule
\endhead
\bottomrule
\endfoot
trees, the order (\cref{def:trees,def:order}) & \lean{Semantics.Tree}, \lean{Semantics.Tree.Le}, \lean{Semantics.polarity} \\
trees over $\BF$ & \lean{Semantics.PrefixSet}, \lean{Semantics.IsRegularTree}, \lean{Semantics.IsFiniteTree} \\
classes of trees & \lean{Full.Logic.TreeClass}, \lean{Full.Models.RegularTop}, \lean{Full.Models.FiniteTop} \\
sentences (\cref{sec:prelim-sentences}) & \lean{Logic.QuantifierFreeFormula}, \lean{Logic.Sigma1}, \lean{Full.Logic.BooleanSentence}, \lean{HoldsIn} \\
flat constraints, entailment (\cref{def:flat,def:disj}) & \lean{Constraints.LitT}, \lean{Constraints.SatT}, \lean{Full.DisjEntailsIn}, \lean{Full.HasCountermodelIn} \\
cones, atoms, rules (\cref{def:cones,tab:clauses}) & \lean{Full.Closure.leaf}, \lean{Config}, \lean{Move}, \lean{Clause}, \lean{Fact}, \lean{Neg}, \lean{IsSolution}, \lean{leastClosure} \\
seeds, exclusions (\cref{sec:readers-modes}) & \lean{Full.seed}, \lean{Full.exclusion} \\
runs (\cref{def:runs}) & \lean{Automata.upperRun}, \lean{Automata.lowerRun}, \lean{upperStep}, \lean{lowerStep} \\
readers, channels, holes (\cref{def:reader,def:hole}) & \lean{Readers.StructuralReader}, \lean{Readers.Table.channel}, \lean{StructuralReader.Eligible}, \lean{IsHole}, \lean{firstRejection} \\
the reader of a mode (\cref{def:compile}) & \lean{Full.Reduction.compile} \\
domination, fresh, immortal (\cref{def:domination,def:fresh,def:immortal}) & \lean{Readers.Dominates}, \lean{FreshDominates}, \lean{IsMinimalFresh}, \lean{StructuralReader.Immortal} \\
the cover (\cref{def:cover-hole}) & \lean{Holes.cover}, \lean{Holes.coverProjection} \\
markers, anchors (\cref{def:markers}) & \lean{Holes.Fresh.level}, \lean{marker}, \lean{active}, \lean{contextAnchors}, \lean{anchors}, \lean{topLevel} \\
parses, trie, block deletion (\cref{def:trie,def:block-deletion}) & \lean{upperOccurrences}, \lean{upperParse}, \lean{trieNodes}, \lean{Edge}, \lean{Edge.TraversedBy}, \lean{label}, \lean{deletionMargin}, \lean{BlockDeletion} \\
rays, ray system (\cref{def:ray,def:ray-system}) & \lean{Words.Ray}, \lean{Ray.onset}, \lean{Ray.period}, \lean{Holes.Tail.raySystem} \\
suffix nodes, closed factors (\cref{def:node}) & \lean{suffixDictionary}, \lean{suffixNode}, \lean{NodeClosedAt}, \lean{NoSuffixOccurrenceEndsBetween}, \lean{pumpFactor} \\
suffix chains & \lean{Words.suffixReturnGap}, \lean{Words.SuffixChain}, \lean{ChainBoundFrom} \\
scales (\cref{def:cover,app:tails-levels}) & \lean{LevelBounds}, \lean{levelStep}, \lean{levelBounds}, \lean{CoordinateCover}, \lean{LevelCoverProperty} \\
forbidden overlaps (\cref{def:overlaps}) & \lean{PrefixOverlap}, \lean{PeriodicTailOverlap}, \lean{ReversePrefixOverlap} \\
pieces, admissible deletions (\cref{def:admissible}) & \lean{Holes.Periodic.Piece}, \lean{FollowerEdge}, \lean{AdmissibleDeletion}, \lean{editedTuple} \\
deletion of positions & \lean{Words.deletePositions}, \lean{retainedCount}, \lean{greatestPreimage} \\
\end{longtable}}

{\small
\begin{longtable}{@{}>{\raggedright\arraybackslash}p{0.30\linewidth}>{\raggedright\arraybackslash\footnotesize}p{0.66\linewidth}@{}}
\caption{Results.}\label{tab:lean-thms}\\
\toprule
Paper & \small Lean \\
\midrule
\endfirsthead
\toprule
Paper & \small Lean \\
\midrule
\endhead
\bottomrule
\endfoot
\cref{thm:main} & \lean{decideBooleanSentence}, \lean{decideBooleanSentence_correct}, \lean{booleanSentence_regular_iff}, \lean{decideBooleanSentenceTopFree}, \lean{booleanSentenceTopFree_regular_iff} \\
\cref{lem:clause,prop:sentences} & \lean{Logic.QuantifierFreeFormula.holds_iff_dnf}, \lean{Logic.ConjunctiveClause.realizable_iff_not_disjEntailsIn}, \lean{Logic.decideBooleanSentence}, \lean{Logic.BooleanSentence.holdsIn_iff_of_disjEntailsIn_iff} \\
the table of \cref{sec:prelim-constraints} & \lean{Full.Logic.holdsIn_satisfiabilitySentence_iff}, \lean{Full.Logic.holdsIn_entailmentSentence_iff} \\
\cref{ex:unary}, \eqref{eq:arrow-disj} & \lean{Examples}: \lean{unary_comparable}, \lean{not_unary_comparable_single}, \lean{binary_not_comparable}, \lean{arrow_least}, \lean{not_arrow_least_single} \\
\cref{thm:cones,lem:clause-form} & \lean{Full.Closure.treeLe_iff_leaf}, \lean{leaf_decode}, \lean{sat_iff_isSolution}, \lean{sat_iff_leastClosure} \\
\cref{lem:duality,lem:witness} & \lean{Full.Closure.moves_dual}, \lean{Full.failures_iff_exists_leaf_witnesses} \\
\cref{thm:realization} & \lean{Full.Reduction.SeedConsistency}, \lean{hasCountermodelIn_iff_exists_seedConsistency}, \lean{seedConsistency_iff_terminalRootConditions} \\
\cref{lem:runs-det,lem:cuts}, \cref{lem:lc} & \lean{Automata.lowest_point}, \lean{Full.Reduction.reach_iff_event}, \lean{Full.Closure.leastClosure_membership} \\
\cref{thm:reduction} & \lean{Full.Reduction.hole_iff_terminalRootConditions}, \lean{not_disjEntails_iff_hole} \\
\cref{lem:hole-dominates,lem:distinct} & \lean{StructuralReader.IsHole.of_dominates}, \lean{runPrefix_ne_of_mortal}, \lean{runPrefix_ne_of_le_firstRejection} \\
\cref{lem:cover,thm:small-hole} & \lean{Holes.cover_fresh}, \lean{cover_project_of_dom}, \lean{Holes.exists_small_hole}, \lean{smallHoleBound}, \lean{Holes.decideHole} \\
\cref{thm:fresh,thm:fresh-bound} & \lean{Holes.Fresh.totalLength_le_freshTupleBound}, \lean{exists_freshDominates_totalLength_le}, \lean{freshTupleBound_le_double_exp2} \\
\cref{lem:markers,lem:depth-span}, \cref{lem:width} & \lean{marker_eq_of_infix}, \lean{marker_occurrences_eq_of_nested}, \lean{contextAnchors_marker_succ}, \lean{topLevel_le}, \lean{card_anchors_Ioc_marker_le}, \lean{card_prefixOccurrenceCuts_le} \\
\cref{lem:parse-local,lem:traverse} & \lean{relativeAnchors_eq}, \lean{BlockDeletion.traversedBy_of_text_eq} \\
\cref{thm:block-deletion} & \lean{contextAnchors_apply}, \lean{BlockDeletion.transfer}, \lean{BlockDeletion.channel}, \lean{BlockDeletion.freshDominates} \\
\cref{lem:labels-local} & \lean{label_local}, \lean{exists_blockDeletion} \\
\cref{thm:tails,thm:tails-bound} & \lean{Holes.Tail.exists_hole_rejectedTailLength_le}, \lean{tailLengthBound}, \lean{tailLengthBound_le_double_exp2} \\
\cref{lem:ray-prefixes} & \lean{ray_prefix_hole}, \lean{Arithmetic.DifferenceSystem.exists_long_solution} \\
\cref{lem:pump,lem:closed-free} & \lean{comparison_of_pumpFactor}, \lean{pumpFactor_eligible}, \lean{exists_closed_factor_free_of_suffix_occurrences}, \lean{Words.chain_stop_from} \\
\cref{lem:levels-mono,lem:chains} & \lean{levelCoverProperty_succ}, \lean{exists_stable_level}, \lean{levelCoverProperty_chainBound}, \lean{levelCoverProperty_chainBoundFrom}, \lean{Words.chainBound_mortal} \\
\cref{lem:terminal,lem:candidates} & \lean{exists_terminal_ray}, \lean{exists_many_rayCandidates}, \lean{exists_compatible_rayCandidate} \\
\cref{lem:compatible} & \lean{card_prefixOverlap_le}, \lean{card_periodicTailOverlap_le}, \lean{card_reversePrefixOverlap_le}, \lean{not_mem_closed_walk_of_noForbiddenRayOverlap} \\
\cref{prop:run-cover} & \lean{exists_rayFamily}, \lean{exists_hole_with_run_cover} \\
\cref{lem:admissible} & \lean{Holes.Periodic.editedTuple_dominates}, \lean{editedTuple_firstRejection}, \lean{exists_preserved_piece_of_nonconjugate_root} \\
\cref{lem:cap}, \cref{lem:simple-path,lem:pass} & \lean{cap_exists}, \lean{simplePath_le_pow_split}, \lean{Arithmetic.simplePath_le_pow_prefix}, \lean{exists_hole_compressed_step} \\
\cref{thm:elementary} & \lean{smallHoleBound_le} \\
\cref{thm:decide-arbitrary,lem:sat-decidable} & \lean{decideDisjEntails}, \lean{Full.Closure.decideSat} \\
\cref{thm:decide-regular} & \lean{Full.Reduction.disjEntails_regular_iff}, \lean{SeedConsistency.regular_realizes}, \lean{Full.Closure.regular_of_leaf} \\
\cref{thm:decide-finite,lem:pinning} & \lean{Full.Reduction.disjEntails_finite_iff}, \lean{Full.Models.finiteBoundPinning}, \lean{observation_bound}, \lean{Constraints.Ranked.Extension.satFin_iff} \\
\cref{lem:guard,thm:decide-bf} & \lean{Semantics.Tree.isTopFree_of_le_of_le_node}, \lean{Entailment.disjEntailsIn_iff_guarded} and its variants for regular and finite trees \\
\end{longtable}}

\subsection{Differences between the Paper and the Development}
\label{app:lean-differences}

The paper and the development follow the same route. They differ in the
following points of representation and arrangement.

\begin{enumerate}
\item \emph{Representation.} Argument positions are numbered from $0$, the
  variance is a Boolean function with \emph{true} for contravariant
  positions, and the signs $\sgp,\sgm$ are the Booleans \emph{false} and
  \emph{true}. Constraints are lists of literals over an arbitrary finite
  type of variables with decidable equality. The alphabet and the state set
  of a reader are arbitrary finite types, and the rejected states and
  admitted pairs are Boolean functions.
\item \emph{Runs.} The development defines the upper and lower runs for
  covariant constraints and applies them to the \emph{signed translation} of
  a constraint with variance, a covariant constraint on two copies of each
  variable; the paper defines them directly on atoms. The reader of a mode
  tests rejection against the set of signed variables from which a variable
  of $\NegV$ is reachable by moves that keep the path empty, which is
  equivalent on upper runs because they are closed under such moves.
\item \emph{Realization.} The development states \cref{thm:realization}
  with a condition for each seed (no forbidden atom, the opposite not in the
  least closure, no derivation of its own opposite) and conditions for pairs
  of different seeds, and proves the equivalence with the form used by the
  readers; the paper states the second form directly.
\item \emph{Labels.} A block deletion requires equal labels in every
  traversing occurrence, as in \cref{def:block-deletion}; the pigeonhole
  uses the pair of the marker texts and the map $g_d$ of
  \cref{lem:labels-local} as a joint label.
\item \emph{Closed factors.} Closed factors for escaping coordinates are
  found among cuts that are $B+1$ apart, which makes the factor longer than
  its window and yields the hypothesis of the count for the overlap
  \ref{it:O2}. An alternative is to add to the suffix node, for each
  terminal ray, the longest suffix of the text that is a factor of the
  periodic part of the ray, with a capped length; the development and the
  paper use the spaced cuts.
\item \emph{Chains.} The chain budget of a fixed word is a bound on the
  number of members of a suffix chain that starts at a long window
  (\cref{lem:chains}). An alternative bounds the number of maximal stretches
  of a chain with a constant gap and the length of each stretch; counting
  members is simpler and gives the same form of the recursion.
\item \emph{State periods.} The development uses the period $d\cdot N!$ of
  the states inside a piece of period $d$, as the paper does; the least
  common multiple of $1,\dots,N$ in place of $N!$ would also do.
\item \emph{Scale functions.} The development packages the inequalities
  \ref{it:L1}--\ref{it:L5}, and a few monotonicity statements, in a
  structure, proves the tail bound from the structure, and then
  instantiates it with the recursion of \cref{app:tails-levels}.
\item \emph{Finite trees.} \Cref{lem:pinning} is proved in the development
  as in \cref{app:classes-finite}: bounds by finite trees are compiled into
  one-sided constructor literals on fresh variables, such an extension
  creates no cycle clash, and a finite solution of the signed translation is
  made symmetric under the exchange of the two copies by a median. The
  development proves the characterization of finite satisfiability by cycle
  clashes itself, for constraints with variable inequalities, whose closure
  contains these inequalities; the paper cites it from \cite{Zhao26}, where
  it is stated without variable inequalities, and removes them first.
\item \emph{Degenerate cases.} The development covers arity $0$, the case
  of no queries, empty words and readers with a single state without
  separate treatment, with truncated subtraction in the numerical
  definitions; the paper treats $k=0$ separately (\cref{sec:tails}) and
  assumes $k\ge1$ and $N\ge2$ in \cref{app:tails,app:compression}, which
  holds for every reader with at least one start and a hole.
\item \emph{Complexity.} The development proves the numerical bound of
  \cref{thm:elementary}. It does not contain \cref{cor:complexity},
  \cref{prop:hardness} or the construction of \cref{app:complexity-long}.
\end{enumerate}
